\lysh{38790}{4}

\begin{alist}
\item Η εξίσωση για τη θερμοκρασία $T$ που αποτυπώνει αυτό που παρατήρησε η Ελένη είναι η:
$|T - 23| = 2$

\item $|T - 23| \le 5 \Leftrightarrow -5 \le T - 23 \le 5 \Leftrightarrow -5 + 23 \le T \le 5 + 23 \Leftrightarrow 18 \le T \le 28$.
Το εύρος των «άνετων» θερμοκρασιών $T$ που επιτρέπει το σύστημα θέρμανσης του σπιτιού της Ελένης είναι από $18^\circ\text{C}$ έως και $28^\circ\text{C}$.

\item Όταν είναι ενεργό το σύστημα ψύξης:
\begin{align*}
18 \le T_1(x) \le 28 \ &\Leftrightarrow \\
18 \le \dfrac{1}{3}x + 17 \le 28 \ &\Leftrightarrow \\
18 - 17 \le \dfrac{1}{3}x \le 28 - 17 \ &\Leftrightarrow \\
1 \le \dfrac{1}{3}x \le 11 \ &\Leftrightarrow
\end{align*}
$3 \le x \le 33$ (1)
Ο τύπος $T_1(x)$ όμως ισχύει για θερμοκρασίες $28 < x \le 40$ (2).
Κάνουμε συναλήθευση των δύο ανισώσεων (1) και (2) και προκύπτει $28 < x \le 33$. Άρα το εύρος των εξωτερικών θερμοκρασιών $x$, για τις οποίες η εσωτερική θερμοκρασία του σπιτιού διακυμαίνεται εντός του άνετου εύρους, όταν λειτουργεί το σύστημα ψύξης, είναι μεγαλύτερο από $28^\circ\text{C}$ έως και $33^\circ\text{C}$.

\item Η πρόταση που περιγράφει λεκτικά τον τύπο $T$ της γραφικής παράστασης είναι η \textbf{1}.
\end{alist}
